Peralatan yang digunakan dalam penelitian ini disiapkan oleh peneliti dan disesuaikan dengan kebutuhan perekaman data bahasa isyarat.  
Seluruh peralatan dirancang untuk mendukung proses pengambilan data yang sederhana, portabel, dan tidak mengganggu kenyamanan partisipan.  

Daftar alat dan bahan yang digunakan disajikan pada Tabel~\ref{tab:alat-bahan}.  

\begin{table}[H]
\centering
\caption{Daftar Alat dan Bahan Penelitian}
\label{tab:alat-bahan}
\begin{tabular}{p{5cm} p{6cm}}
\hline
\textbf{Alat/Bahan} & \textbf{Keterangan} \\
\hline
Laptop atau komputer & Digunakan untuk perekaman, pemrosesan data, dan pelatihan model \\
Kamera webcam Logitech 1080p & Digunakan sebagai perangkat utama perekaman video \\
Kamera tambahan & Digunakan apabila diperlukan untuk sudut pandang alternatif \\
Tripod kamera & Menjaga kestabilan posisi kamera selama perekaman \\
Perangkat lunak MediaPipe & Digunakan untuk ekstraksi \textit{landmark} tubuh, tangan, dan wajah \\
Lingkungan pemrograman Python & Digunakan untuk pengolahan data dan pemodelan \\
\hline
\end{tabular}
\end{table}

Peneliti bertanggung jawab atas penyediaan dan pengoperasian seluruh peralatan penelitian.  
Apabila pihak PUSBISINDO memiliki ketentuan, rekomendasi, atau fasilitas pendukung tambahan, peneliti siap menyesuaikan pelaksanaan penelitian sesuai dengan arahan yang diberikan.  